\section{Разработка программного комплекса}

\subsection{Архитектура и потоки данных}

Программный комплекс спроектирован с использованием модульной архитектуры, нацеленной на разделение ответственности (Separation of Concerns). Это обеспечивает высокую гибкость, простоту тестирования и дальнейшего расширения. Проект логически разделен на несколько пакетов, взаимодействие которых формирует основной поток данных в приложении.

\begin{figure}[h]
\centering
\includegraphics[width=0.8\textwidth]{architecture.png} % Предполагается наличие схемы
\caption{Схема архитектуры приложения}
\label{fig:arch}
\end{figure}

\textbf{Основные пакеты и их роли:}
\begin{itemize}
    \item \texttt{cmd/analyzer}: Точка входа. Отвечает за инициализацию приложения Fyne, создание главного окна и сборку основного контейнера с вкладками. Не содержит бизнес-логики.
    \item \texttt{pkg/gui}: Реализация пользовательского интерфейса. Пакет содержит конструкторы для каждой панели (вкладки). Панели запрашивают данные у модулей \texttt{profiling}, \texttt{benchmark}, \texttt{memory} и \texttt{concurrency}, а также управляют запуском тестов и отображением результатов. Взаимодействие с логикой происходит асинхронно через горутины.
    \item \texttt{pkg/profiling}: Модуль сбора системных метрик. Предоставляет функции для получения данных о ЦП, ОЗУ и т.д. Является «поставщиком» данных для панели мониторинга.
    \item \texttt{pkg/benchmark}, \texttt{pkg/memory}, \texttt{pkg/concurrency}: Модули, содержащие логику самих синтетических тестов. Они не зависят от GUI и могут быть использованы отдельно. Каждый тест представляет собой функцию, возвращающую измеряемую величину.
\end{itemize}

\textbf{Поток данных при мониторинге:} Панель мониторинга (\texttt{DashboardPanel}) запускает фоновую горутину, которая с определенным интервалом вызывает функции из пакета \texttt{profiling}. Полученные данные (строки) используются для обновления текстовых меток (\texttt{widget.Label}) в интерфейсе.

\textbf{Поток данных при бенчмаркинге:} Пользователь нажимает кнопку «Запустить тесты» на одной из панелей. Обработчик нажатия запускает горутину, которая итерируется по списку тестов. Для каждого теста вызывается универсальный \texttt{benchmark.TestRunner}, который, в свою очередь, многократно выполняет функцию конкретного теста. Результаты (статистика) возвращаются в горутину панели и используются для обновления соответствующей метки результата.

\subsection{Реализация модуля бенчмаркинга}

Это ядро всего комплекса. Рассмотрим реализацию ключевых тестов.

\subsubsection{Измерение задержки доступа к памяти (Latency)}

Для измерения задержки доступа к кэшам L1/L2/L3 и ОЗУ используется метод \textbf{«погони за указателями» (pointer chasing)}. Его суть — создать длинную цепочку зависимостей по данным, чтобы исключить влияние аппаратных предсказателей (prefetcher) и измерить «чистое» время доступа.

\textbf{Алгоритм:}
1. Создается большой массив байтов, размер которого заведомо превышает размер тестируемого уровня кэша.
2. Массив заполняется индексами так, чтобы сформировать псевдослучайную последовательность обхода. Например, \texttt{array[i] = j, array[j] = k, ...}. Это создает длинный цикл, проходящий через все элементы массива в непредсказуемом для процессора порядке.
3. Запускается цикл, который многократно проходит по этой цепочке указателей: \texttt{p = array[p]}. Каждая итерация зависит от результата предыдущей, что не позволяет процессору выполнять несколько итераций параллельно (out-of-order execution) и заставляет его ожидать данные из памяти.
4. Измеряется общее время выполнения и делится на количество операций, чтобы получить среднюю задержку на одну операцию доступа.

Подбирая размер массива, можно «попасть» в определенный уровень кэша. Если массив помещается в L1, мы измеряем задержку L1. Если он больше L1, но меньше L2, мы измеряем задержку L2 (с учетом промахов в L1), и так далее.

\begin{lstlisting}[language=Go, caption={Фрагмент кода для измерения задержки}]
func runLatencyTest(size int, iterations int) float64 {
    data := make([]int, size)
    // ... создание псевдослучайной цепочки индексов ...

    start := time.Now()
    p := 0
    for i := 0; i < iterations; i++ {
        p = data[p] // Pointer chasing
    }
    duration := time.Since(start)

    return float64(duration.Nanoseconds()) / float64(iterations)
}
\end{lstlisting}

\subsubsection{Измерение пропускной способности памяти (Bandwidth)}

Здесь задача иная — загрузить шину памяти по максимуму. Для этого выполняются последовательные операции чтения или записи на большом непрерывном блоке памяти.

\begin{itemize}
    \item \textbf{Тест на запись:} В цикле всему массиву присваивается одно и то же значение. Компилятор Go достаточно умен, чтобы не оптимизировать этот цикл.
    \item \textbf{Тест на чтение:} В цикле значения из массива читаются в одну и ту же локальную переменную. Это предотвращает оптимизацию, при которой чтение могло бы быть полностью удалено.
\end{itemize}

Пропускная способность вычисляется как \( (\text{Размер данных} \times \text{Количество итераций}) / \text{Время выполнения} \).

\subsubsection{Тесты механизмов параллелизма}

\begin{itemize}
    \item \textbf{Накладные расходы на горутины:} Создается и сразу же завершается большое количество горутин. Измеряется общее время, что позволяет оценить, насколько «дешево» обходится создание одной горутины.
    \item \textbf{Блокировка мьютекса:} Несколько горутин в цикле пытаются захватить один и тот же мьютекс (\texttt{sync.Mutex}), инкрементировать счетчик и освободить мьютекс. Тест показывает, насколько высока цена синхронизации при высоком уровне конкуренции.
    \item \textbf{Пропускная способность каналов:} Одна горутина (писатель) отправляет данные в канал, а другая (читатель) их получает. Измеряется, сколько таких пересылок можно совершить в секунду.
\end{itemize}

\subsection{Анализ и интерпретация результатов}

После проведения тестов приложение отображает средние значения. Важно правильно интерпретировать эти цифры. 

\textbf{Пример анализа результатов теста задержки памяти:}
Предположим, получены следующие результаты:
\begin{itemize}
    \item L1 Cache Latency: 0.9 ns/op
    \item L2 Cache Latency: 3.5 ns/op
    \item L3 Cache Latency: 15 ns/op
    \item RAM Latency: 70 ns/op
\end{itemize}
Из этих данных видна четкая иерархия памяти. Доступ к L1 почти в 4 раза быстрее, чем к L2, и в ~80 раз быстрее, чем к ОЗУ. Это наглядно демонстрирует важность написания программ, которые эффективно используют кэши (принцип локальности данных).

\textbf{Пример анализа пропускной способности ОЗУ:}
\begin{itemize}
    \item Read Bandwidth: 25 ГБ/с
    \item Write Bandwidth: 22 ГБ/с
\end{itemize}
Эти цифры можно сравнить с теоретической пиковой пропускной способностью, заявленной производителем памяти (например, для DDR4-3200 это 25.6 ГБ/с на один канал). Расхождение показывает, насколько эффективно система использует доступную полосу пропускания.

\textbf{Пример анализа тестов параллелизма:}
Сравнение производительности мьютекса и атомарных операций для инкремента счетчика может показать, что на современных процессорах атомарные операции (\texttt{atomic.AddInt64}) в десятки раз быстрее, чем использование мьютекса для защиты одной переменной. Это дает разработчику понимание, какой примитив синхронизации выбрать в конкретной ситуации.
